\section{Additional Analysis Protocols}
\label{app:additional-analysis}

\subsection{Matched-Compute Component Ablation}
\label{app:ablation-protocol}

All variants use the same student initialization, optimizer, epochs, seeds, and scored anchor--column budget. Batch-local KD normalizes over co-occurring examples; random corpus support draws a uniformly sampled support of matched size; one-step RIPPLE uses only direct mutual-neighbor rows; diffusion adds broader targets; and full RIPPLE adds cross-anchor row supervision. An injection-only control uses the same graph traversals and candidate replacement with $\lambda_{\mathrm{row}}=0$. A degree-preserving random-graph control rewires adjacency before applying the same targets, sampler, and loss construction.

The primary locality comparison reports downstream average, full-corpus teacher--student JS divergence, and retained teacher mass. The extended table additionally reports teacher-neighbor recall and student off-support mass. This separation prevents low restricted-support KL from being interpreted as global agreement.

\subsection{Diffusion and Graph Width}

One-hop and diffusion variants are evaluated on teacher--student error in fixed teacher-rank bands and on held-out diffusion profiles at scales broader than those directly supervised. Graph-width sweeps keep candidate budget and optimization fixed. Because increasing $k$ trivially changes available support, every width is evaluated on the same fixed rank bands as well as on full-corpus divergence.

The degree-preserving random-graph control tests whether gains follow from teacher geometry rather than generic graph regularization. For each rewired graph, degree sequence, target scales, candidate quota, and update count are held fixed.

\subsection{Candidate Budget and Gradient Fidelity}

Candidate-budget sweeps report retained teacher mass, selected-support distortion $\bar\epsilon$, missed mass $\bar\delta$, their sum, full-corpus distortion $\mathcal E$, full-corpus JS, downstream average, step time, and peak memory. For small corpus slices where full scoring is feasible, gradient fidelity is measured by the cosine similarity between the sparse-candidate gradient and the full-corpus relational gradient under the same anchors and student checkpoint.

\subsection{Cross-Anchor Row Diagnostics}

Row-supervision comparisons report distinct selected rows, row hit rate, the effective number of available columns, transition-row KL, and disagreement on graph rows shared by multiple anchors. The injection-only control separates the row loss from changes in candidate composition. All diagnostics use the same anchor set and pooled candidates across compared variants.
